\section{Implementation}
\label{sec:implementation}

We implement \texttt{MaybeUncertain<T>} in the Rust crate
\texttt{deep\_causality\_uncertain}, which is part of the DeepCausality
project~\cite{hansen_deepcausality}. The implementation reuses the
\texttt{UncertainNodeContent} enum, the \texttt{SequentialSampler},
the \texttt{sprt\_eval} module, and the global sample cache already used
for $\Utype\langle T \rangle$; the only new code is a struct, four
constructors, five operator overloads, the \texttt{sample} method, two
presence predicates, and the SPRT-gated lifting operator, comprising
roughly two hundred lines of Rust across two files. This section
presents the relevant API surface and the design choices behind it.

\subsection{Type and trait bounds}
\label{sec:impl:type}

The struct declaration is two-fielded and generic over a type parameter
bounded by the crate's \texttt{ProbabilisticType} trait:

\begin{lstlisting}
pub struct MaybeUncertain<T: ProbabilisticType> {
    is_present: Uncertain<bool>,
    value:      Uncertain<T>,
}
\end{lstlisting}

The trait \texttt{ProbabilisticType} extends \texttt{IntoSampledValue},
\texttt{FromSampledValue}, \texttt{Clone}, \texttt{Send}, \texttt{Sync},
and \texttt{'static}; it provides a \texttt{default\_value} constant for
the placeholder used in \texttt{always\_none}. The crate provides
implementations for \texttt{f64} and \texttt{bool}.

Both struct fields are private. Callers construct values exclusively
through the four constructors of~\Cref{sec:type:ctors}, and they observe
values through \texttt{sample}, the presence predicates, and the gating
operator. The library exposes two type aliases for the supported
specialisations:

\begin{lstlisting}
pub type MaybeUncertainF64  = MaybeUncertain<f64>;
pub type MaybeUncertainBool = MaybeUncertain<bool>;
\end{lstlisting}

\subsection{Constructors}
\label{sec:impl:ctors}

The four constructors map directly onto~\Cref{tab:constructors}. For the
\texttt{f64} specialisation:

\begin{lstlisting}
pub fn from_value(value: f64) -> Self {
    Self {
        is_present: Uncertain::<bool>::point(true),
        value:      Uncertain::<f64>::point(value),
    }
}

pub fn from_uncertain(value: Uncertain<f64>) -> Self {
    Self { is_present: Uncertain::<bool>::point(true), value }
}

pub fn always_none() -> Self {
    Self {
        is_present: Uncertain::<bool>::point(false),
        value:      Uncertain::<f64>::point(0.0),
    }
}

pub fn from_bernoulli_and_uncertain(
    prob_some:          f64,
    present_value_dist: Uncertain<f64>,
) -> Self {
    Self {
        is_present: Uncertain::bernoulli(prob_some),
        value:      present_value_dist,
    }
}
\end{lstlisting}

The placeholder $0.0$ in \texttt{always\_none} is never observed at
runtime because the sampler short-circuits when the presence draw is
false. The \texttt{bool} specialisation is structurally identical except
that the placeholder is \texttt{false} instead of $0.0$.

\subsection{Sampling}
\label{sec:impl:sampling}

The two-stage sampling procedure from~\Cref{alg:sample} is implemented
as a single method that delegates to the underlying \texttt{sample}
methods on the two component types:

\begin{lstlisting}
pub fn sample(&self) -> Result<Option<f64>, UncertainError> {
    if self.is_present.sample()? {
        Ok(Some(self.value.sample()?))
    } else {
        Ok(None)
    }
}
\end{lstlisting}

Each \texttt{Uncertain::sample} call draws a fresh random sample index from
the thread-local generator and queries the global sample cache keyed by
$(\text{node id}, \text{sample index})$. Two consequences follow. First,
repeated \texttt{sample} calls on the same \texttt{MaybeUncertain} value
yield independent draws, since the indices differ between calls. Second,
within a single sampling operation, the underlying graph evaluates shared
sub-nodes exactly once because the cache lookup succeeds for any leaf
reached through more than one path.

The SPRT loop uses a stricter variant, \texttt{sample\_with\_index}, which
takes a caller-supplied index. The loop in
\texttt{sprt\_eval::evaluate\_hypothesis} sweeps a contiguous range of
indices starting from a base, drawing one bool per index and accumulating
log-likelihood ratios. This is the only place in the crate that exercises
the deterministic-index path; ordinary user code uses \texttt{sample},
which is non-deterministic by design.

\subsection{Operator overloads}
\label{sec:impl:ops}

The five operators on \texttt{MaybeUncertain<f64>} are implemented through
Rust's \texttt{std::ops} traits. Every binary operator follows the same
pattern; the \texttt{Add} implementation is representative:

\begin{lstlisting}
impl Add for MaybeUncertain<f64> {
    type Output = Self;
    fn add(self, rhs: Self) -> Self::Output {
        Self {
            is_present: self.is_present & rhs.is_present,
            value:      self.value + rhs.value,
        }
    }
}
\end{lstlisting}

Two distinct things happen on each call. The value field uses
\texttt{Uncertain<f64>}'s overloaded \texttt{+}, which appends an
\texttt{ArithmeticOp} node to the computation graph rooted at
\texttt{value}. The presence field uses \texttt{BitAnd} on
\texttt{Uncertain<bool>}, which appends a \texttt{LogicalOp::And} node to
the computation graph rooted at \texttt{is\_present}. Both ops consume the
operands by value, taking advantage of Rust's affine-type discipline; no
hidden cloning is performed.

The \texttt{Sub}, \texttt{Mul}, and \texttt{Div} implementations are
structurally identical, differing only in the arithmetic operator applied
to the value branch. The unary \texttt{Neg} implementation leaves the
presence axis untouched:

\begin{lstlisting}
impl Neg for MaybeUncertain<f64> {
    type Output = Self;
    fn neg(self) -> Self::Output {
        Self { is_present: self.is_present, value: -self.value }
    }
}
\end{lstlisting}

\subsection{Presence predicates and the gating operator}
\label{sec:impl:lift}

The two presence predicates are one-line shims that clone the
\texttt{is\_present} field and optionally negate it:

\begin{lstlisting}
pub fn is_some(&self) -> Uncertain<bool> {  self.is_present.clone()   }
pub fn is_none(&self) -> Uncertain<bool> { !self.is_present.clone()   }
\end{lstlisting}

The \texttt{!} in \texttt{is\_none} is the \texttt{Not} impl on
\texttt{Uncertain<bool>} and constructs a \texttt{LogicalOp::Not} node
rather than computing a Boolean. Both predicates therefore return new
\texttt{Uncertain<bool>} values that can be sampled, composed with other
Boolean operators, or passed to any of the existing SPRT-backed methods on
\texttt{Uncertain<bool>}.

The gating operator \texttt{lift\_to\_uncertain} is the only method that
introduces new control flow specific to \texttt{MaybeUncertain<T>}:

\begin{lstlisting}
pub fn lift_to_uncertain(
    &self,
    threshold_prob_some: f64,
    confidence_level:    f64,
    epsilon:             f64,
    max_samples:         usize,
) -> Result<Uncertain<f64>, UncertainError> {
    let is_present = self.is_present.to_bool(
        threshold_prob_some,
        confidence_level,
        epsilon,
        max_samples,
    )?;
    if is_present {
        Ok(self.value.clone())
    } else {
        Err(UncertainError::PresenceError(
            "Insufficient evidence for presence".to_string(),
        ))
    }
}
\end{lstlisting}

The method delegates the statistical decision to
\texttt{Uncertain<bool>::to\_bool}, which in turn calls the SPRT
implementation described in~\Cref{sec:semantics:sprt}. If the SPRT
accepts the presence hypothesis the method returns a clone of the value
graph; otherwise it returns a \texttt{PresenceError}. The error is a
distinct \texttt{UncertainError} variant so that callers can handle
gate-failure separately from sampling or graph-construction errors.

\subsection{Scope and current limitations}
\label{sec:impl:scope}

We state honestly what the current implementation does not provide.

\paragraph{Specialisations.} Only \texttt{f64} and \texttt{bool} are
supported. Generalising to other numeric types would require extending the
inherent-impl blocks; the type itself is generic.

\paragraph{Boolean algebra.} The \texttt{MaybeUncertain<bool>}
specialisation does not implement logical operators. Boolean expressions
over multiple presence-uncertain values require the caller to project the
components, operate on \texttt{Uncertain<bool>} directly, and reconstruct.

\paragraph{Closed-form moments.} The crate provides no closed-form
expressions for the mean, variance, or any other moment of a
\texttt{MaybeUncertain<T>} value. All numerical summaries are estimated by
sampling.

\paragraph{Joint presence-value distribution.} The two components are
modelled as independent random variables. There is no constructor that
takes a joint distribution; missing-not-at-random workloads are out of
scope (\Cref{sec:discussion}).

\paragraph{Conditional and comparison operators.} The crate does not
currently lift conditional expressions
(\texttt{Uncertain::conditional}) or scalar comparison operators onto
\texttt{MaybeUncertain<T>}. Callers needing them work on the underlying
components.
